\pdfbookmark[0]{\IfLanguageName{ngerman}{Zusammenfassung}{Abstract}}{abstract}
\begin{center} 
    \huge \IfLanguageName{ngerman}{Zusammenfassung}{Abstract}
\end{center}

Synthetic communities (SynComs) are a promising approach to increase plant productivity and protection. Furthermore, they provide a controlled framework to study plant-microbiome associations. Still, the impact of assembly strategy on community behaviour is not fully understood. In this thesis I reconstructed, curated, analysed and compared genome-scale metabolic models of two SynComs stemming from the barley rhizosphere: a core community assembled from co-occurring strains, and a trait-optimised community composed of strains previously shown to benefit plant growth. The reconstructed and curated models meet current community quality standards, providing a reliable basis for community-level analyses. Both communities rely heavily on cross-feeding to meet their metabolic requirements but differ strongly in the structure of these exchanges: the core community shows more symmetric, evenly distributed metabolite exchanges, while the trait-optimised community's exchanges are sparser, of larger magnitude and more asymmetric. Further, the core community is substantially more robust to single-strain removal, whereas the trait-optimised community's growth fluctuates depending on which member is excluded. Drop-out experiments identified a keystone strain \textit{Bosea sp.} F3-2, whose removal strongly reduced core community growth, and its addition to the trait-optimised community increased community growth. Together, these findings demonstrate the profound impact of assembly strategy on SynCom behaviour, indicating that co-occurrence fosters stability indicators that are not ensured by trait-based assembly priorities. This has direct implications for the rational design of stable, plant-beneficial SynComs.